\chapter{Suggested Final Rotation Artifact}
\label{ch:final-rotation-artifact}

\textit{A strong rotation artifact is not just a collection of simulations. It is a compact research package: a reproduction, a derivation, a mesoscopic extension, a microscopic comparison, a phase diagram, and a clear statement of what remains unresolved.}

% ============================================================
\section{A Clean Reproduction Notebook}
\label{sec:clean-reproduction-notebook}
% ============================================================

The first artifact should be a single clean reproduction notebook or script-backed notebook that rebuilds the E-clustered baseline from a configuration file. Its purpose is not exploratory plotting. Its purpose is to answer: can we reproduce the minimal phenomena reliably?

The notebook should have five outputs:
%
\begin{enumerate}[leftmargin=*]
  \item a raster plot with cluster labels visible;
  \item smoothed excitatory cluster rates and inhibitory population rate;
  \item active-cluster count $n_{\mathrm{act}}(t)$;
  \item Fano factor and autocorrelation summaries;
  \item a table of parameter values and random seeds.
\end{enumerate}

The core state variable should match the theory:
%
\begin{equation}
  \widehat r_k(t)
  =
  \frac{1}{N_k}
  \sum_{i\in \mathcal{C}_k}(\kappa*s_i)(t).
  \label{eq:artifact-cluster-rate}
\end{equation}
%
If the notebook makes claims about cluster up-states, those claims should be made from $\widehat r_k(t)$ or a documented thresholded version of it.

The reproduction notebook should keep all parameters in one visible configuration block. A future reader should be able to change cluster strength, cluster size, or random seed without editing analysis code.

% ============================================================
\section{A Derivation Note}
\label{sec:derivation-note}
% ============================================================

The second artifact should be a derivation note that connects the microscopic and mesoscopic descriptions. It does not need to be long, but it should be explicit about the retained variables and discarded degrees of freedom.

The minimal derivation structure is:
%
\begin{enumerate}[leftmargin=*]
  \item define microscopic spike trains $s_i(t)$ and cluster activities $A_k^a(t)$;
  \item define filtered rates $r_k^a(t)$;
  \item write the effective inputs $u_k^a(t)$ in terms of coupling matrices;
  \item introduce transfer functions $\Phi_a$;
  \item add finite-size noise with the intended $N_k^{-1/2}$ scaling;
  \item state what assumptions are approximations rather than consequences.
\end{enumerate}

The central equation should be the E/I mesoscopic model:
%
\begin{equation}
  \tau_a\dot r_k^a
  =
  -r_k^a
  +
  \Phi_a(u_k^a)
  +
  \frac{\sigma_a(r_k^a)}{\sqrt{N_k^a}}\xi_k^a(t),
  \qquad
  a\in\{E,I\}.
  \label{eq:artifact-ei-mesoscopic-equation}
\end{equation}
%
Each term should be interpreted: relaxation, expected population response, and intrinsic finite-size fluctuation.

The derivation note should also include the fast-inhibition reduction from \cref{eq:ei-effective-coupling}. That reduction is the shortest mathematical explanation of how inhibitory clustering changes effective competition among excitatory clusters.

% ============================================================
\section{E-Only Mesoscopic Model Results}
\label{sec:e-only-mesoscopic-results}
% ============================================================

The E-only mesoscopic model is the control theory model. It should reproduce the qualitative behavior of the E-clustered microscopic baseline before any E/I extension is presented.

The result set should include:
%
\begin{enumerate}[leftmargin=*]
  \item fixed-point branches as a function of cluster strength or drive;
  \item local stability eigenvalues for low and high cluster states;
  \item stochastic trajectories showing metastable switching;
  \item dwell-time distributions and size scaling;
  \item comparison to microscopic cluster-rate statistics.
\end{enumerate}

A useful minimal conclusion has the form
%
\begin{equation}
  g_{EE}\uparrow
  \Rightarrow
  \text{bistable cluster fixed points}
  \Rightarrow
  \text{finite-size switching when } n_E<\infty .
  \label{eq:e-only-result-chain}
\end{equation}
%
The first arrow should be supported by deterministic fixed points. The second should be supported by stochastic simulations and cluster-size scaling.

This E-only result prevents the project from over-attributing effects to inhibitory clustering. If the E-only mesoscopic model already explains a statistic, the E/I extension should be asked what it improves: robustness, local balance, spike irregularity, parameter range, or match to a specific task model.

% ============================================================
\section{E/I Mesoscopic Extension}
\label{sec:ei-mesoscopic-extension-artifact}
% ============================================================

The E/I extension artifact should introduce paired inhibitory clusters as the smallest additional mechanism. It should make one explicit comparison against the E-only model while holding the baseline as fixed as possible.

The central comparison is:
%
\begin{equation}
  \theta_{\mathrm{E-only}}
  \longrightarrow
  \theta_{\mathrm{E/I}}
  =
  \theta_{\mathrm{E-only}}
  \cup
  \{g_{IE},g_{EI},g_{II},N_k^I,\tau_I,\Phi_I\}.
  \label{eq:ei-extension-parameter-addition}
\end{equation}
%
This notation makes clear which parameters are new. It also discourages retuning every old parameter until the desired result appears.

The E/I result should report:
%
\begin{enumerate}[leftmargin=*]
  \item local E/I tracking $\operatorname{corr}_t(r_k^E,r_k^I)$;
  \item balance residuals $B_k(t)$;
  \item stability region in $(g_{EE},g_{EI})$ or $(g_{IE},g_{EI})$;
  \item whether within-state spike irregularity changes;
  \item whether metastable dwell times become more robust to parameter perturbations.
\end{enumerate}

The strongest possible conclusion would not be that E/I clustering merely creates metastability. E-only clustering already can do that. The stronger conclusion would be that E/I clustering preserves metastability while improving local balance or expanding the region where microscopic spiking remains irregular and stable.

% ============================================================
\section{Microscopic vs Mesoscopic Comparison}
\label{sec:microscopic-vs-mesoscopic-comparison-artifact}
% ============================================================

The comparison should be built around matched variables. For each microscopic simulation, compute the empirical mesoscopic state
%
\begin{equation}
  \widehat{\vec x}(t)
  =
  \left(
    \widehat r_1^E(t),\ldots,\widehat r_Q^E(t),
    \widehat r_1^I(t),\ldots,\widehat r_Q^I(t)
  \right).
  \label{eq:empirical-mesoscopic-state-vector}
\end{equation}
%
Then compare it to $\vec x(t)$ from the mesoscopic equations using statistics that do not require spike-by-spike identity.

The comparison table should include:
%
\begin{longtable}{@{}p{0.24\textwidth}p{0.34\textwidth}p{0.32\textwidth}@{}}
\toprule
\textbf{Observable} & \textbf{Microscopic estimate} & \textbf{Mesoscopic estimate} \\
\midrule
Mean cluster rate & Smoothed spikes averaged over neurons. & Time average of $r_k^a(t)$. \\
State occupancy & Thresholded $\widehat r_k^a(t)$. & Thresholded $r_k^a(t)$. \\
Dwell time & Consecutive active intervals. & Same detection rule. \\
Fano factor & Spike counts in windows. & Counts sampled from rate, or integrated activity approximation. \\
Autocorrelation & $\widehat r_k^a(t)$ or $n_{\mathrm{act}}(t)$. & $r_k^a(t)$ or active-cluster count. \\
Transition path & Transition-triggered average of smoothed spikes. & Transition-triggered average of rates. \\
\bottomrule
\end{longtable}

The same detection thresholds and smoothing windows should be used wherever possible. If they differ, the reason should be stated. Otherwise a good match can be manufactured by analysis choices rather than by model quality.

% ============================================================
\section{Phase Diagram}
\label{sec:phase-diagram-artifact}
% ============================================================

The phase diagram is the summary artifact that turns model behavior into a research result. It should include at least one E-only diagram and one E/I diagram with aligned axes where possible.

A minimal phase-diagram package contains:
%
\begin{enumerate}[leftmargin=*]
  \item categorical regime map with predefined labels;
  \item scalar overlay: leading eigenvalue, $\log\E[L]$, autocorrelation time, or Fano factor;
  \item representative trajectories for selected grid points;
  \item seed and simulation-duration table;
  \item explicit classification rules.
\end{enumerate}

The diagram should distinguish deterministic and stochastic classifications. For example, a point may have multiple stable deterministic fixed points but show no switches in a finite stochastic run because barriers are too high. Another point may show irregular deterministic trajectories with positive Lyapunov exponent. These are not the same regime.

The phase diagram should answer one rotation-level question: does E/I clustering change the mechanism, or does it change the robustness of the same mechanism?

% ============================================================
\section{Summary of Open Questions}
\label{sec:summary-open-questions-artifact}
% ============================================================

The final artifact should end with open questions that are precise enough to become next experiments. Useful open questions include:
%
\begin{enumerate}[leftmargin=*]
  \item Which parameters most strongly shape the mesoscopic transfer functions?
  \item Does the E/I extension require paired inhibitory clusters, or would structured $I\to E$ feedback alone suffice?
  \item Does switching vanish under fixed-$Q$, large-cluster scaling?
  \item Are any irregular mesoscopic regimes deterministic with positive Lyapunov exponent?
  \item Which variability statistic is best matched: Fano factor, autocorrelation, dwell time, or state occupancy?
  \item Which claims depend on exact parameter values from the source papers and therefore need citation-level verification?
\end{enumerate}

Each open question should include a proposed test. An open question without a test is too vague for a rotation handoff.

% ============================================================
\section{Possible Next-Paper Contribution}
\label{sec:possible-next-paper-contribution}
% ============================================================

A plausible next-paper contribution would be a mesoscopic theory of E/I clustered attractor networks that connects three levels: microscopic spiking architecture, E/I mesoscopic equations, and a phase diagram with scaling predictions. In shorthand,
%
\begin{equation}
  \text{micro}
  \longrightarrow
  \text{E/I meso}
  \longrightarrow
  \text{phase diagram}.
  \label{eq:next-paper-three-levels}
\end{equation}
%
The contribution is strongest if it makes a prediction that is hard to see from simulation alone.

Examples of defensible contribution statements are:
%
\begin{enumerate}[leftmargin=*]
  \item E/I clustering expands the parameter region where attractor-like cluster states coexist with balanced irregular spiking.
  \item The apparent slow variability in a specified regime is finite-size-driven because switching disappears under fixed-$Q$ cluster-size scaling.
  \item A different specified regime is deterministic mesoscopic irregularity because switching persists with intrinsic noise removed and the largest Lyapunov exponent is positive.
  \item The fast-inhibition effective coupling matrix predicts which E/I architectures favor single-cluster versus multi-cluster active states.
\end{enumerate}

The final paper-like story should avoid overclaiming. It can state that the mesoscopic model preserves some observables and fails on others. That is scientifically more useful than presenting a polished model with no failure modes.

\begin{chaptersummary}

\begin{center}
\renewcommand{\arraystretch}{1.45}
\begin{tabularx}{\textwidth}{@{}l X X@{}}
  \toprule
  \textbf{Artifact} & \textbf{Purpose} & \textbf{Main evidence} \\
  \midrule
  Reproduction notebook & Establish baseline phenomena & Raster, cluster rates, Fano factor, lifetimes \\
  Derivation note & Connect microscopic and mesoscopic variables & Equations, assumptions, noise scaling \\
  E-only mesoscopic result & Control model for clustering & Fixed points, stability, switching \\
  E/I extension & Test paired inhibition mechanism & Balance, robustness, phase-region changes \\
  Micro/meso comparison & Validate retained variables & Matched observables and scaling \\
  Phase diagram & Turn simulations into theory & Regime map with classification rules \\
  Open questions & Define next experiments & Testable uncertainties \\
  \bottomrule
\end{tabularx}
\end{center}

\end{chaptersummary}

\begin{conceptquestions}
  \item Why should the reproduction notebook expose all parameters in one configuration block?
  \item What would count as evidence that the E/I extension improves robustness rather than merely retuning the E-only model?
  \item Which microscopic observable is most directly matched to the mesoscopic variable $r_k^E(t)$?
  \item What is one next-paper claim that would require a Lyapunov-exponent calculation?
\end{conceptquestions}
